%! Author = javierdesant, morteega, dagalle, jorge
%! Date = 28/09/2026

\documentclass[11pt,a4paper]{article}

% --- Definición de los autores ---
\newcommand{\autoruno}{Javier de Santiago Soto}
\newcommand{\autordos}{Marcos Ortega Jiménez}
\newcommand{\autortres}{Daniel }    %FIXME
\newcommand{\autorcuatro}{Jorge }   %FIXME

% --- Definición de reglas y variables ---
\newenvironment{tight-list}{
    \begin{itemize}
    \setlength{\itemsep}{0pt}
    \setlength{\parskip}{0pt}
}{\end{itemize}}

% --- Paquetes de idioma y tipografía ---
\usepackage[utf8]{inputenc}
\usepackage[T1]{fontenc}
\usepackage[spanish,es-tabla]{babel}
\usepackage{microtype}
\usepackage{lmodern}

% --- Geometría y márgenes ---
\usepackage[top=2.5cm,bottom=2.5cm,left=2.5cm,right=2.5cm]{geometry}

% --- Paquetes matemáticos y símbolos ---
\usepackage{amsmath}
\usepackage{amsfonts}
\usepackage{amssymb}

% --- Tablas y figuras ---
\usepackage{graphicx}
\usepackage{booktabs}
\usepackage{tabularx}
\usepackage{multirow}
\usepackage{float}
\usepackage{caption}
\usepackage{subcaption}

% --- Formato de código y comandos Junos ---
\usepackage{listings}
\usepackage{xcolor}

\definecolor{codegray}{rgb}{0.5,0.5,0.5}
\definecolor{codebg}{rgb}{0.96,0.96,0.96}
\definecolor{keywordblue}{rgb}{0.13,0.13,0.7}

\lstdefinestyle{junosstyle}{
    backgroundcolor=\color{codebg},
    basicstyle=\ttfamily\small,
    breakatwhitespace=false,
    breaklines=true,
    captionpos=b,
    keepspaces=true,
    numbers=left,
    numbersep=8pt,
    numberstyle=\tiny\color{codegray},
    showspaces=false,
    showstringspaces=false,
    showtabs=false,
    tabsize=2,
    frame=single,
    framerule=0.5pt,
    rulecolor=\color{codegray}
}
\lstset{style=junosstyle}

% --- Enlaces y referencias cruzadas ---
\usepackage{hyperref}
\hypersetup{
    colorlinks=true,
    linkcolor=black,
    citecolor=blue,
    urlcolor=blue,
    pdfauthor={\autoruno, \autordos, \autortres, \autorcuatro},
    pdftitle={Memoria\_Practica\_1}
}

% --- Datos del documento ---
\title{\textbf{Memoria Técnica: Diseño de un Centro de Datos \textit{Bare-Metal}}}
\author{\autoruno \and \autordos \and \autortres \and \autorcuatro}
\date{Curso 2026--2027}

\begin{document}    % TODO: update both abstracts when finished

    \maketitle

    \begin{abstract}
        Esta memoria técnica presenta el diseño integral de la infraestructura de red
        para un centro de datos \textit{bare-metal} compuesto por \textbf{6.822 servidores}
        en \textbf{360 racks}. A partir de una matriz de tráfico agregado de \textbf{17,34~Tbps}
        marcadamente vertical (84\,\% norte-sur), se propone una arquitectura Fat-Tree
        modular no bloqueante adaptada al perfil de distribución de contenidos.

        El diseño integra redundancia activa en las capas de acceso y agregación, cortafuegos
        perimetrales y un esquema de direccionamiento IPv4 jerárquico optimizado para
        sumarización. La solución garantiza una fiabilidad del \textbf{99,958\,\%}
        (indisponibilidad inferior a 3,7~h/año) y valida su viabilidad económica con un
        coste de cableado en la lista de materiales (\textit{BoM}) de \textbf{92.340~€}
        (13,54~€/servidor). La conectividad y políticas de seguridad se validan mediante un
        banco de pruebas en laboratorio sobre equipamiento Juniper Networks.
    \end{abstract}

    \renewcommand{\abstractname}{Abstract}
    \begin{abstract}
        This technical report presents the comprehensive network infrastructure design
        for a \textit{bare-metal} data center hosting \textbf{6,822 servers} distributed
        across \textbf{360 racks}. Based on an aggregate traffic matrix of \textbf{17.34~Tbps}
        with a predominantly vertical flow pattern (84\,\% north-south), a modular non-blocking
        Fat-Tree architecture tailored for content delivery workloads is proposed.

        The architecture incorporates active redundancy across the access and aggregation layers,
        perimeter firewalls, and a hierarchical IPv4 addressing scheme optimized for route
        summarization. The proposed solution delivers a system availability of \textbf{99.958\,\%}
        (annual downtime under 3.7~hours) and confirms economic feasibility with a bill of materials
        (\textit{BoM}) cabling expenditure of \textbf{€92,340} (€13.54 per server). End-to-end
        connectivity and security policies are empirically validated through a laboratory testbed
        deployed on Juniper Networks hardware.
    \end{abstract}

    \newpage
    \tableofcontents
    \newpage
    \listoffigures
    \listoftables
    \newpage

% --- Secciones Modulares ---
    \section{Introducción y Alcance del Proyecto}
\label{sec:introduccion}

El objetivo de este proyecto es diseñar la arquitectura de red de un centro de
datos que dispone únicamente de infraestructura \textit{bare metal}, a partir
de los requisitos de tráfico y de servidores proporcionados en el enunciado.
La red debe dar servicio a \textbf{6.822 servidores} distribuidos en
\textbf{360 racks} de cómputo, maximizando la utilidad del sistema al menor coste
posible y cumpliendo un objetivo de fiabilidad del \textbf{99,958\,\%}.

\subsection{Objetivos del Proyecto}
\label{subsec:objetivos}

Este objetivo general se concreta en los siguientes objetivos específicos:

\begin{tight-list}
    \setlength{\itemsep}{0pt}
    \setlength{\topsep}{2pt}
    \setlength{\parsep}{0pt}
    \item \textbf{Caracterizar el tráfico.}
    \item \textbf{Dimensionar las capas de acceso y agregación.}
    \item \textbf{Aplicar la redundancia necesaria.}
    \item \textbf{Dimensionar el \textit{core} y ajustar la sobresuscripción.}
    \item \textbf{Proteger la red mediante cortafuegos.}
    \item \textbf{Definir el direccionamiento y el enrutamiento.}
    \item \textbf{Evaluar el coste.}
    \item \textbf{Validar el diseño en laboratorio.}
\end{tight-list}

Todas las decisiones de diseño se justifican a lo largo de la memoria,
explicando su motivación y las alternativas consideradas para maximizar el rendimiento.

\subsection{Limitaciones y Alcance}
\label{subsec:limitaciones}

Para focalizar los recursos en los componentes de mayor impacto y eficiencia en la red, se establecen las siguientes delimitaciones:

\begin{tight-list}
    \setlength{\itemsep}{0pt}
    \setlength{\topsep}{2pt}
    \setlength{\parsep}{0pt}
    \item Los servidores físicos, los conmutadores ToR y el cableado interno intra-rack se consideran preexistentes e integrados en el bastidor de cómputo.
    \item Quedan fuera del alcance la ingeniería de instalaciones (suministro eléctrico, refrigeración y espacio físico) y la gestión del software de aplicación o licencias superiores.
\end{tight-list}

    \section{Análisis de Tráfico y Dimensionamiento de la Infraestructura}
\label{sec:analisis_trafico}

\subsection{Metodología de Procesamiento de Datos}
\label{subsec:metodologia_trafico}
% TODO: Explicar el preprocesamiento y extracción sobre las muestras de tráfico cada 10 min (grupo_4.xlsx).

\subsection{Caracterización Estadística del Tráfico}
\label{subsec:estadistica_trafico}
% TODO: Incluir métricas de tráfico medio, pico, percentiles (p95/p99) y márgenes de crecimiento.

\subsection{Dimensionamiento de Capacidad y Enlaces}
\label{subsec:dimensionamiento_capacidad}
% TODO: Detallar el cálculo de número de servidores, densidad de puertos ToR, agregación y enlaces de core.

    \section{Diseño de la Arquitectura de Red}
\label{sec:arquitectura_red}

\subsection{Topología Fat-Tree (Clos 3-Tier)}
\label{subsec:topologia_fattree}
% TODO: Describir la estructura multicapa: Acceso (Leaf), Agregación (Spine) y Núcleo (Core).

\subsection{Análisis de Sobresuscripción y Bisección}
\label{subsec:sobresuscripcion}
% TODO: Presentar los ratios de sobresuscripción entre niveles y análisis de probabilidad de congestión.

\subsection{Selección y Roles de Equipamiento Juniper}
\label{subsec:roles_equipamiento}
% TODO: Justificar el rol de cada dispositivo (MX80, SRX320/340, QFX5100).

    \section{Estudio y Modelado Matemático de Fiabilidad y Disponibilidad}
\label{sec:fiabilidad_disponibilidad}

\subsection{Modelado de Fiabilidad y Diagramas de Bloques (RBD)}
\label{subsec:rbd_model}
% TODO: Presentar los diagramas de bloques de fiabilidad para la topología con dual-homing.

\subsection{Formulación Analítica Serie-Paralelo}
\label{subsec:formulacion_matematica}
% TODO: Desarrollar las ecuaciones analíticas de disponibilidad basadas en MTBF y MTTR.

\subsection{Evaluación de Disponibilidad Global \texorpdfstring{($A \ge 99{,}999\%$)}{(A >= 99.999\%)}}
\label{subsec:evaluacion_disponibilidad}
% TODO: Presentar los cálculos de disponibilidad alcanzada y análisis de eliminación de SPOF.

    \section{Plan Detallado de Direccionamiento y Numeración IP}
\label{sec:plan_direccionamiento}

\subsection{Esquema Jerárquico de Direccionamiento IPv4}
\label{subsec:esquema_ipv4}
% TODO: Definir la partición RFC 1918 por capas (Core, Agregación, Acceso, Servidores, Gestión).

\subsection{Segmentación y Asignación de VLANs 802.1Q}
\label{subsec:segmentacion_vlans}
% TODO: Tabla de mapeo de VLANs, identificadores, nombres y rangos asignados.

\subsection{Matriz de Interconexiones Punto a Punto}
\label{subsec:matriz_p2p}
% TODO: Detallar subredes /30 o /31 para los enlaces entre routers, switches y firewalls.

    \section{Tablas de Enrutamiento y Políticas de Seguridad}
\label{sec:enrutamiento_seguridad}

\subsection{Estrategia de Enrutamiento y Tablas de Rutas}
\label{subsec:estrategia_enrutamiento}
% TODO: Especificar rutas estáticas/dinámicas, saltos siguientes y resúmenes de red en MX80 y QFX5100.

\subsection{Zonas de Seguridad en Cortafuegos SRX}
\label{subsec:zonas_seguridad}
% TODO: Definir zonas (trust, untrust, dmz, mgmt) e interfaces asociadas en SRX320/340.

\subsection{Políticas de Tránsito y Control de Tráfico de Host}
\label{subsec:politicas_seguridad}
% TODO: Especificar reglas de filtrado de paquetes y servicios permitidos hacia el host (SSH, ICMP).

    \section{Configuración Detallada para Equipos Juniper (Junos CLI)}
\label{sec:configuracion_junos}

\subsection{Configuración de Routers de Core Juniper MX80}
\label{subsec:cfg_mx80}
% TODO: Incluir bloques de comandos Junos (inicialización, direccionamiento, modo packet-based, enrutamiento).

\subsection{Configuración de Cortafuegos Juniper SRX320 / SRX340}
\label{subsec:cfg_srx}
% TODO: Incluir bloques de comandos Junos (zonas de seguridad, políticas, host-inbound-traffic).

\subsection{Configuración de Conmutadores Juniper QFX5100}
\label{subsec:cfg_qfx}
% TODO: Incluir bloques de comandos Junos (VLANs 802.1Q, interfaces access y trunk).

    \section{Presupuesto, Estimación de Costes y Lista de Materiales (BOM)}
\label{sec:presupuesto_bom}

\subsection{Lista de Materiales de Hardware y Óptica (BOM)}
\label{subsec:bom_hardware}
% TODO: Tabla con desglose de equipos (MX80, SRX320/340, QFX5100), transceptores y cables.

\subsection{Estimación de Costes de Inversión (CapEx) y Operación (OpEx)}
\label{subsec:costes_capex_opex}
% TODO: Detallar costes unitarios, totales y mantenimiento de la infraestructura.

    \section{Plan de Pruebas y Validación en Laboratorio}
\label{sec:plan_pruebas}

\subsection{Batería de Pruebas de Conectividad y Enrutamiento}
\label{subsec:pruebas_conectividad}
% TODO: Definir pruebas de ping ICMP, traceroute y verificación de saltos entre pods.

\subsection{Validación de Políticas de Seguridad y Aislamiento}
\label{subsec:pruebas_seguridad}
% TODO: Comprobación de bloqueo de tráfico no autorizado y acceso a zonas restringidas.

\subsection{Pruebas de Tolerancia a Fallos y Reconvergencia}
\label{subsec:pruebas_tolerancia_fallos}
% TODO: Simulación de caída de enlaces/switches y validación de dual-homing.

    \section{Conclusiones y Trabajo Futuro}
\label{sec:conclusiones}

\subsection{Conclusiones del Diseño}
\label{subsec:conclusiones_diseno}
% TODO: Resumir los logros del dimensionamiento, fiabilidad y viabilidad del diseño propuesto.

\subsection{Líneas de Trabajo Futuro}
\label{subsec:trabajo_futuro}
% TODO: Identificar posibles ampliaciones (BGP EVPN/VXLAN, automatización NetDevOps, telemetría).


\end{document}